% ECONOMICS_PROBLEM: {"id":"C73-11","title":"Characterizing universal convergence of fictitious play","jel_code":"C73","source_id":"OP-0136"}
% FIRST_POSTED: 2026-10-09
% ATTEMPT_STATUS: unresolved
% ENTRY_KIND: research_attempt
\documentclass[11pt]{article}
\usepackage{amsmath,amssymb,amsthm,geometry,hyperref}
\geometry{margin=1in}
\newtheorem{theorem}{Theorem}
\newtheorem{proposition}[theorem]{Proposition}
\newtheorem{lemma}[theorem]{Lemma}
\newtheorem{corollary}[theorem]{Corollary}
\theoremstyle{definition}
\newtheorem{example}[theorem]{Example}
\newtheorem{remark}[theorem]{Remark}
\newcommand{\R}{\mathbb R}
\newcommand{\N}{\mathbb N}
\DeclareMathOperator{\NE}{NE}
\DeclareMathOperator{\cav}{cav}
\DeclareMathOperator{\supp}{supp}
\title{Characterizing universal convergence of fictitious play\\ GPT 6 Astra Ultra}
\author{Research attempt}
\date{9 October 2026}
\begin{document}
\maketitle

\section*{Target and scope}
The target is a payoff-array characterization of simultaneous fictitious play that approaches the Nash set for every initial empirical profile and every legal tie choice. The update is
\[
x_i^{t+1}=\frac{t x_i^t+e_{a_i^{t+1}}}{t+1},\qquad
 a_i^{t+1}\in\arg\max_a u_i(a,x_{-i}^t).
\]
The universal characterization remains unresolved. This attempt gives a quantitative, checkable sufficient condition and proves directly that it is not necessary. Both are elementary restricted results, not proposed new classes beyond the established fictitious-play literature. The survey \cite{survey} is the source of the broader classification agenda. No theorem about convergence of continuous-time or asynchronous play is substituted for this discrete rule.

\section*{A quantitative elimination argument}
Assume $|u_i(s)|\le M$ for all pure profiles, with $M>0$. A sequential strict-dominance deletion certificate is a finite sequence of restricted action sets
\[
 A^0_i=S_i,\quad A^j_i\subseteq A^{j-1}_i\quad(1\le j\le m),
\]
where at step $j$ exactly one action $b_j$ of player $i_j$ is removed. There is a probability vector $q_j$ supported on $A^j_{i_j}$ and a number $\gamma_j>0$ such that
\[
 u_{i_j}(q_j,s_{-i_j})-u_{i_j}(b_j,s_{-i_j})\ge\gamma_j
 \quad\text{for every }s_{-i_j}\in\prod_{k\ne i_j}A^{j-1}_k.
\]
The certificate is called complete if every $A_i^m$ is a singleton.

\begin{theorem}[Explicit extinction of choices and convergence rate]
For every finite game, every such certificate, every initial $x^1$, and every legal fictitious-play tie sequence, define recursively
\[
 B_0=0,\qquad
 T_j=\max\left\{1,T_{j-1},\left\lceil\frac{8M B_{j-1}}{\gamma_j}\right\rceil\right\},
 \qquad B_j=B_{j-1}+T_j,
\]
where $T_0=1$. Then $b_j$ is never selected at an update based on time $t\ge T_j$, and
\[
 x_{i_j,b_j}^t\le\frac{T_j}{t}\quad(t\ge T_j).
\]
If the certificate is complete, its remaining pure profile $s^*$ is the unique Nash equilibrium, and for $t\ge T_m$,
\[
 \|x^t-e_{s^*}\|_2\le\frac{2B_m}{t}.
\]
\end{theorem}
\begin{proof}
Induct on $j$. At time $t\ge T_{j-1}$, the sum of the probabilities of all previously removed actions, summed over all players, is at most $B_{j-1}/t$ by the induction hypothesis. Under independent opponent mixing, the probability $\theta$ that at least one opponent uses a previously removed action is at most this sum, by the union bound. The payoff difference between $q_j$ and $b_j$ is at least $\gamma_j$ when all opponents use surviving actions and at least $-2M$ otherwise. Its expectation is therefore at least
\[
 (1-\theta)\gamma_j-2M\theta
 \ge\gamma_j-4M B_{j-1}/t,
\]
since $\gamma_j\le 2M$. For $t\ge T_j$ this is at least $\gamma_j/2$; when $B_{j-1}=0$ it is at least $\gamma_j$ directly. A pure best response cannot be strictly worse than the mixed action $q_j$. Thus $b_j$ cannot be chosen at any of these updates. The averaging rule gives $t x_{i_j,b_j}^t=T_j x_{i_j,b_j}^{T_j}\le T_j$, proving the inductive claim.

For uniqueness of equilibrium, first note that a strictly dominated action receives zero probability in any equilibrium of its restricted game: the mixture dominating it has strictly larger payoff against the equilibrium opponents, whereas every action assigned positive probability must maximize payoff. Inducting over deletions shows that any equilibrium of the original game must use only the final actions.

The final profile is indeed an equilibrium. Fix player $i$ and its final opponents. At each deletion of one of $i$'s actions, its payoff is strictly below a convex combination of payoffs of actions still available. Consequently the maximum payoff among $i$'s available actions does not fall when that action is removed. Deletions of other players leave the fixed final opponents unchanged. Hence $i$'s last action attains the original maximum. This proves equilibrium existence here directly. Finally, the mass outside the final actions, summed over players, is at most $B_m/t$. The sum of the absolute coordinate differences from the pure profile is twice that mass, and the Euclidean norm is at most the sum of absolute coordinate differences.
\end{proof}

The time constants can be large, but their role is explicit: beliefs retain initial and historical mass, so an action dominated only after deletions does not disappear from beliefs in finite time. It stops being selected and then decays as $1/t$. This distinction is essential for the stated update.

\section*{A full tie-robust calculation outside that certificate}
\begin{proposition}[The binary matching coordination game]
Let both players have actions $0,1$ and receive payoff $1$ when their actions match and $0$ otherwise. Every simultaneous fictitious-play trajectory from every pair of initial mixed strategies converges to a Nash equilibrium, for arbitrary time-dependent legal tie choices. This game has no strict-dominance deletion certificate removing even its first action.
\end{proposition}
\begin{proof}
Write $p_t=x_{1,1}^t$, $q_t=x_{2,1}^t$ and
\[
 D_t=t(2p_t-1),\qquad E_t=t(2q_t-1).
\]
Then $D_{t+1}=D_t+s_t$ and $E_{t+1}=E_t+r_t$, where $s_t=1$ if $E_t>0$, $s_t=-1$ if $E_t<0$, and $s_t\in\{-1,1\}$ if $E_t=0$. The analogous rule determines $r_t$ from $D_t$.

If both $D_t,E_t$ are strictly positive at some time, they increase by one at every subsequent time. Thus $(p_t,q_t)\to(1,1)$. If both are strictly negative, they decrease by one forever and $(p_t,q_t)\to(0,0)$.

Suppose neither event ever occurs. We show both coordinates stay bounded. Let $C=\max\{|D_1|,|E_1|,2\}$. If $D_t>1$, then $E_t\le0$. The case $E_t<0$ decreases $D_t$ by one. If $E_t=0$, then $E_{t+1}=1$; to avoid the strictly positive quadrant one must have $D_{t+1}\le0$, which is impossible when $D_t>1$. Thus $D_t>1$ implies that its next step decreases it. Symmetrically, $D_t<-1$ implies that its next step increases it. When $|D_t|\le1$, its next value has absolute value at most two. The same reasoning applies to $E_t$. Induction gives $|D_t|,|E_t|\le C$, and hence $(p_t,q_t)\to(1/2,1/2)$.

These three profiles are exactly the Nash equilibria: each player's strict best reply is the action favored by the opponent's probability, with indifference only at $1/2$. Neither pure action can be strictly dominated by a mixture, because it gives payoff one against the identical opposing pure action, the maximal possible payoff. Thus no first deletion is possible.
\end{proof}

\section*{Verification and remaining gap}
The preceding proof handles boundary initial beliefs and ties exactly. For example, an antisymmetric pair $D_t=-E_t$ may keep crossing zero; the bounded-coordinate case makes this converge to the mixed equilibrium rather than incorrectly declaring eventual pure play. Strict dominance is therefore only sufficient, while a convergence classification must also include games with multiple equilibria. No necessary-and-sufficient payoff condition for arbitrary finite games is proved here.

\begin{thebibliography}{9}
\bibitem{survey} H. Li, W. Huang, Z. Duan, D. H. Mguni, K. Shao, J. Wang, and X. Deng, \emph{A survey on algorithms for Nash equilibria in finite normal-form games}, arXiv:2312.11063 (2023), learning-dynamics discussion and concluding open problems. \url{https://arxiv.org/pdf/2312.11063}. Checked against the primary preprint on 9 October 2026. The restricted proofs above are supplied in full and do not depend on transferring a convergence theorem from a different updating convention.
\end{thebibliography}

\end{document}
